\documentclass[a4paper,12pt]{article}

% ─── Кодировка и язык ───
\usepackage[utf8]{inputenc}
\usepackage[T2A]{fontenc}
\usepackage[russian]{babel}

% ─── Поля ───
\usepackage[left=2.5cm, right=1.5cm, top=2cm, bottom=2cm]{geometry}

% ─── Пакеты ───
\usepackage{graphicx}
\usepackage{booktabs}
\usepackage{longtable}
\usepackage{hyperref}
\usepackage{listings}
\usepackage{xcolor}
\usepackage{caption}
\usepackage{enumitem}
\usepackage{float}
\usepackage{tabularx}

% ─── Настройка листингов ───
\lstset{
  basicstyle=\ttfamily\small,
  backgroundcolor=\color{gray!10},
  frame=single,
  breaklines=true,
  breakatwhitespace=true,
  showstringspaces=false,
  numbers=left,
  numberstyle=\tiny\color{gray},
  keywordstyle=\color{blue!70},
  commentstyle=\color{green!50!black},
  stringstyle=\color{red!60},
  language=Python,
  inputencoding=utf8,
  extendedchars=true,
  literate={а}{{\cyra}}1 {б}{{\cyrb}}1 {в}{{\cyrv}}1 {г}{{\cyrg}}1
           {д}{{\cyrd}}1 {е}{{\cyre}}1 {ж}{{\cyrzh}}1 {з}{{\cyrz}}1
           {и}{{\cyri}}1 {й}{{\cyrishrt}}1 {к}{{\cyrk}}1 {л}{{\cyrl}}1
           {м}{{\cyrm}}1 {н}{{\cyrn}}1 {о}{{\cyro}}1 {п}{{\cyrp}}1
           {р}{{\cyrr}}1 {с}{{\cyrs}}1 {т}{{\cyrt}}1 {у}{{\cyru}}1
           {ф}{{\cyrf}}1 {х}{{\cyrh}}1 {ц}{{\cyrc}}1 {ч}{{\cyrch}}1
           {ш}{{\cyrsh}}1 {щ}{{\cyrshch}}1 {ы}{{\cyry}}1 {ь}{{\cyrsftsn}}1
           {э}{{\cyrerev}}1 {ю}{{\cyryu}}1 {я}{{\cyrya}}1
}

\hypersetup{
  colorlinks=true,
  linkcolor=blue!60!black,
  urlcolor=blue!70!black,
  citecolor=green!50!black
}

% ─── Межстрочный интервал ───
\usepackage{setspace}
\onehalfspacing

\begin{document}

% ═══════════════════════════════════════════
% ТИТУЛЬНЫЙ ЛИСТ
% ═══════════════════════════════════════════

\begin{titlepage}
\begin{center}

\textbf{РОССИЙСКИЙ УНИВЕРСИТЕТ ДРУЖБЫ НАРОДОВ\\ИМЕНИ ПАТРИСА ЛУМУМБЫ}\\[0.3cm]
\textbf{Факультет физико-математических и естественных наук}\\[0.3cm]
Направление подготовки: 09.03.03 <<Прикладная информатика>>\\[2cm]

{\Large \textbf{ЗАДАНИЕ 5П}}\\[0.5cm]
{\large \textbf{Прототип системы с документацией}}\\[1cm]

{\large Тема НИР:}\\[0.3cm]
{\large <<Анализ тональности сообщений пользователей социальных сетей\\
и оценка их влияния на динамику цен криптовалют>>}\\[2cm]

\begin{tabular}{ll}
\textbf{Студент:} & Мугари Абдеррахим \\
\textbf{Группа:} & НФИбд-01-22 \\
\textbf{Науч. рук.:} & доц., к.т.н. Молодченков А.И. \\
\textbf{Трек:} & П (прикладной) \\
\end{tabular}

\vfill
Москва, 2026

\end{center}
\end{titlepage}

\tableofcontents
\newpage

% ═══════════════════════════════════════════
% 1. ОПИСАНИЕ ПРОТОТИПА
% ═══════════════════════════════════════════

\section{Описание прототипа}

\subsection{Назначение системы}

Разработанный прототип представляет собой программный комплекс для автоматизированного
анализа тональности сообщений пользователей Twitter, посвящённых криптовалюте Bitcoin,
и оценки статистической взаимосвязи между агрегированными показателями тональности
и динамикой рыночных цен BTC/USD.

\subsection{Ключевой функционал}

Прототип реализует следующие функциональные блоки:

\begin{enumerate}[label=\arabic*.]
  \item \textbf{Загрузка и фильтрация данных} --- чтение датасета Twitter-сообщений (CSV),
        фильтрация по языку (английский), дедупликация по идентификатору твита.
  \item \textbf{Предобработка текста} --- удаление URL-адресов, упоминаний пользователей (@),
        префикса ретвита (RT), приведение к нижнему регистру, токенизация,
        удаление стоп-слов (NLTK).
  \item \textbf{Анализ тональности (VADER)} --- лексиконный метод, использующий
        оригинальный текст с учётом регистра, пунктуации и эмодзи. Генерирует
        составной показатель (compound) в диапазоне $[-1, 1]$.
  \item \textbf{Анализ тональности (FinBERT)} --- трансформерная модель
        ProsusAI/finbert (109.5M параметров), обученная на финансовых текстах.
        Классифицирует каждый твит как positive/neutral/negative.
  \item \textbf{Получение ценовых данных} --- API CoinGecko с fallback-механизмом
        (встроенные исторические данные при недоступности API).
  \item \textbf{Агрегация и объединение} --- ежедневная агрегация показателей тональности
        и объединение с ценовыми данными BTC.
  \item \textbf{Корреляционный анализ} --- лагированные корреляции Пирсона и Спирмена
        (лаги 0--7 дней), тест Грейнджер-каузальности, тест Дики--Фуллера на стационарность.
  \item \textbf{Визуализация} --- автоматическая генерация 7 графиков и 6 таблиц
        в форматах PNG и CSV.
\end{enumerate}

\subsection{Репозиторий}

Код проекта размещён на GitHub:

\begin{center}
\url{https://github.com/iragoum/nlp-crypto-sentiment-pipeline}
\end{center}

% ═══════════════════════════════════════════
% 2. РЕАЛИЗОВАННЫЙ VS ЗАПЛАНИРОВАННЫЙ
% ═══════════════════════════════════════════

\section{Реализованный функционал в сравнении с запланированным}

\begin{table}[H]
\centering
\caption{Сравнение запланированного и реализованного функционала}
\label{tab:planned_vs_done}
\begin{tabularx}{\textwidth}{|p{5.5cm}|c|X|}
\hline
\textbf{Функционал} & \textbf{Статус} & \textbf{Комментарий} \\
\hline
Загрузка и фильтрация Twitter-данных & Реализовано & 664\,918 англоязычных твитов из 847 CSV-файлов \\
\hline
Предобработка текста (очистка, токенизация) & Реализовано & Regex-фильтрация, NLTK-токенизация \\
\hline
VADER-анализ тональности & Реализовано & Составной показатель compound для каждого твита \\
\hline
FinBERT-анализ тональности & Реализовано & ProsusAI/finbert, 3 класса (pos/neu/neg) \\
\hline
Получение цен BTC (CoinGecko API) & Реализовано & С fallback-механизмом для offline-режима \\
\hline
Агрегация на дневном уровне & Реализовано & Объединение sentiment + price по дате \\
\hline
Корреляции Пирсона/Спирмена (лаги 0--7) & Реализовано & 8 лагов, p-value для каждого \\
\hline
Тест Грейнджер-каузальности & Реализовано & F-статистика, p-value, лаги 1--7 \\
\hline
ADF-тест на стационарность & Реализовано & Для VADER и BTC return \\
\hline
Сравнение VADER vs FinBERT & Реализовано & Таблица распределений, матрица согласия \\
\hline
Визуализация (графики) & Реализовано & 7 графиков PNG с подписями осей и легендой \\
\hline
CI/CD (GitHub Actions) & Реализовано & Python 3.10, 3.11, 3.12; pytest + flake8 \\
\hline
torchinfo-вывод архитектуры & Реализовано & 109\,484\,547 параметров \\
\hline
ONNX-экспорт для Netron & Реализовано & Файл .onnx + скриншот в README \\
\hline
Веб-интерфейс & Не реализовано & Запланировано для ВКР \\
\hline
Мультивалютный анализ (ETH, BNB) & Не реализовано & Запланировано для ВКР \\
\hline
\end{tabularx}
\end{table}

Все запланированные для этапа НИР компоненты реализованы в полном объёме.
Веб-интерфейс и мультивалютный анализ запланированы для следующего этапа (ВКР).

% ═══════════════════════════════════════════
% 3. ИНСТРУКЦИЯ ПО ЗАПУСКУ
% ═══════════════════════════════════════════

\section{Инструкция по запуску}

\subsection{Требования}

\begin{itemize}
  \item Python 3.10 или выше
  \item pip (менеджер пакетов Python)
  \item Git
\end{itemize}

\subsection{Установка}

\begin{lstlisting}[language=bash, caption={Установка зависимостей}]
git clone https://github.com/iragoum/nlp-crypto-sentiment-pipeline.git
cd nlp-crypto-sentiment-pipeline
python -m venv venv
source venv/bin/activate     # Linux/Mac
venv\Scripts\activate        # Windows
pip install torch --index-url https://download.pytorch.org/whl/cpu
pip install -r requirements.txt
\end{lstlisting}

\subsection{Запуск}

\begin{lstlisting}[language=bash, caption={Команды запуска}]
# Полный pipeline (VADER + FinBERT)
python main.py

# Только VADER (быстрый режим, ~10 мин)
python main.py --skip-finbert --skip-torchinfo

# Подвыборка для тестирования
python main.py --skip-torchinfo --sample 50000

# Запуск тестов
pytest tests/ -v
\end{lstlisting}

% ═══════════════════════════════════════════
% 4. РЕЗУЛЬТАТЫ ТЕСТИРОВАНИЯ
% ═══════════════════════════════════════════

\section{Результаты тестирования}

\subsection{Модульные тесты (pytest)}

Тестовый набор включает 28 модульных тестов, охватывающих три основных модуля системы.

\begin{table}[H]
\centering
\caption{Результаты модульного тестирования}
\label{tab:test_results}
\begin{tabular}{|l|c|c|}
\hline
\textbf{Модуль} & \textbf{Кол-во тестов} & \textbf{Статус} \\
\hline
\texttt{test\_preprocessor.py} & 10 & 10/10 passed \\
\hline
\texttt{test\_sentiment\_analyzer.py} & 6 & 6/6 passed \\
\hline
\texttt{test\_correlation\_analyzer.py} & 12 & 12/12 passed \\
\hline
\textbf{Итого} & \textbf{28} & \textbf{28/28 passed} \\
\hline
\end{tabular}
\end{table}

\subsection{Покрытие тестами}

\begin{itemize}
  \item \textbf{Preprocessor:} удаление URL, упоминаний, RT-префикса; приведение к нижнему
        регистру; обработка пустых строк и None; токенизация; удаление стоп-слов.
  \item \textbf{Sentiment Analyzer:} корректность колонок VADER; диапазон compound
        $[-1, 1]$; правильность полярности для позитивных и негативных текстов;
        сохранение индекса DataFrame.
  \item \textbf{Correlation Analyzer:} агрегация по дням; корректность лагированных
        корреляций; диапазон коэффициентов; тест Грейнджера; ADF-тест;
        корреляционная матрица.
\end{itemize}

\subsection{CI/CD}

Непрерывная интеграция настроена через GitHub Actions. При каждом push в ветку \texttt{main}
автоматически запускаются:

\begin{enumerate}
  \item Модульные тесты (\texttt{pytest tests/ -v}) на Python 3.10, 3.11, 3.12.
  \item Линтинг кода (\texttt{flake8}).
\end{enumerate}

Статус CI: \textbf{все проверки проходят успешно} (см. бейдж в README репозитория).

\subsection{Метрики работы pipeline}

\begin{table}[H]
\centering
\caption{Метрики выполнения pipeline (50\,000 твитов)}
\label{tab:pipeline_metrics}
\begin{tabular}{|l|c|}
\hline
\textbf{Метрика} & \textbf{Значение} \\
\hline
Общее число твитов в датасете & 664\,918 \\
\hline
Число обработанных твитов (выборка) & 50\,000 \\
\hline
Диапазон дат & 26.07.2022 -- 30.08.2022 \\
\hline
Число уникальных дней & 36 \\
\hline
Средний VADER compound & $+0.169$ \\
\hline
FinBERT: positive / neutral / negative & 3.46\% / 91.89\% / 4.65\% \\
\hline
Согласованность VADER и FinBERT & 47.22\% \\
\hline
Корреляция Пирсона (lag 0) & $r = 0.251$, $p = 0.146$ \\
\hline
Тест Грейнджера (lag 1--7) & Незначим ($p > 0.05$) \\
\hline
ADF-тест (VADER daily mean) & Нестационарен ($p = 0.88$) \\
\hline
ADF-тест (BTC daily return) & Стационарен ($p < 0.001$) \\
\hline
\end{tabular}
\end{table}

% ═══════════════════════════════════════════
% 5. НОВИЗНА ПОДХОДА
% ═══════════════════════════════════════════

\section{Новизна подхода}

Основная новизна прототипа заключается в \textbf{одновременном применении двух
принципиально различных подходов к анализу тональности} и их систематическом
сравнении в контексте криптовалютного рынка:

\begin{enumerate}
  \item \textbf{VADER} (Valence Aware Dictionary and sEntiment Reasoner) ---
        лексиконный метод, учитывающий регистр, пунктуацию, степень интенсивности
        и эмодзи. Не требует обучения, работает на исходном тексте.
  \item \textbf{FinBERT} (ProsusAI/finbert) --- трансформерная модель на базе BERT,
        дообученная на корпусе финансовых текстов. 109.5 млн параметров.
        Работает на очищенном тексте.
\end{enumerate}

Ключевое наблюдение: \textbf{согласованность двух методов составляет лишь 47.22\%},
что свидетельствует о существенном различии в интерпретации тональности
криптовалютных сообщений. FinBERT классифицирует 91.89\% твитов как нейтральные,
тогда как VADER распределяет оценки более равномерно (43.23\% позитивных,
45.59\% нейтральных, 11.18\% негативных).

Это различие открывает перспективу для дальнейшего исследования:
какой из подходов обеспечивает более высокую предсказательную силу
для динамики цен криптовалют.

\end{document}
